\subsection{A.5: a wild counterexample on a regular complete local surface}

\begin{theorem}[Refuted; A.5 even for excellent regular complete local rings]
\label{A:wild-affine-counterexample}
Let $k=\overline{\F}_2$, let $W=W(k)$ be its ring of Witt vectors, and put
\[
 A=W[[t]],\qquad U=\Spec A[1/t].
\]
Then $A$ is an excellent regular henselian local domain of dimension two,
its fraction field has characteristic zero, and
\[
 \cd_2(U)=3>2=\dim A+\delta(k)+\cd_2(k).
\]
In fact $H^3(U,\F_2)\neq0$. Thus the residual-characteristic case of
Conjecture 4 on 31/085--087 is false even with quasi-excellence.
\end{theorem}
\begin{proof}
The ring $W$ is a complete discrete valuation ring with uniformizer $2$
and residue field $k$. Hence $A$ is a complete regular local domain with
regular parameters $(2,t)$. Complete noetherian local rings are excellent.
The residue field is algebraically closed and perfect, so both terms
$\delta(k)$ and $\cd_2(k)$ are zero.

Write
\[
 R=\Spec A[1/2],\quad V=\Spec A[1/(2t)],\quad
 B=k[[t]],\quad E=k((t)),\quad C=\Frac W.
\]
Let $D=(A_{(2)})^h$ be the henselization of the discrete valuation ring
$A_{(2)}$, and let $L=\Frac D$. Its residue field is $E$.
All symbols below are degree-two Kummer symbols with coefficients
$\mu_2^{\otimes2}=\F_2$ on characteristic-zero schemes.

The complement $z=U-V$ is the single closed point $\Spec E$.
Henselian invariance gives $\cd_2(D)=\cd_2(E)\leq1$.
Gabber--Orgogozo gives $d_2(E)=1$ and $\cd_2(L)=2$:
apply their Th\'eor\`eme 1.2 first to $k[[t]]$ and then to the excellent
henselian discrete valuation ring $D$.
By Proposition~\ref{A:wild-generic-open}, both $R$ and $V$ have
$2$-cohomological dimension two.
\'Etale excision for henselization, followed by the localization sequence
for the closed point of $\Spec D$, gives
\[
 H^3_z(U,\F_2)=H^3_z(\Spec D,\F_2)
       \simeq H^2(L,\F_2).
\]
Consequently localization on $U$ gives an exact sequence
\begin{equation}\label{A:wild-cokernel}
 H^2(V,\F_2)\longrightarrow H^2(L,\F_2)
       \longrightarrow H^3(U,\F_2)\longrightarrow0.
\end{equation}
For arbitrary $2$-primary torsion coefficients, the same localization
sequences give vanishing above degree three: the terms on $D$, $L$ and
$V$ vanish above degrees one, two and two respectively. Thus
$\cd_2(U)\leq3$.

We now exhibit an element outside the image in
\eqref{A:wild-cokernel}. The only point of $R-V$ is $(t)$, with residue
field $C$. Purity applies here because $2$ is invertible on $R$.
The localization sequence and $H^3(R,\F_2)=0$ show that
\[
 H^2(R,\F_2)\longrightarrow H^2(V,\F_2)
       \longrightarrow H^1(C,\F_2)\longrightarrow0
\]
is exact. Every element of the final group is represented by some
$u\in C^\times$ and lifted by the symbol $\{t,u\}$, since its residue
at $(t)$ is the Kummer class of $u$. Therefore every class coming from
$H^2(V,\F_2)$ has the form
\begin{equation}\label{A:wild-image-form}
 a+\{t,u\},\qquad a\in H^2(R,\F_2),\quad u\in C^\times.
\end{equation}

Here are the precise ramification-filtration inputs. Let
$U^mH^2(L,\F_2)$ be generated by symbols whose first entry lies in
$1+2^mD$ for $m\geq1$, and put $U^0=H^2$.
Since the absolute ramification index is one, the upper critical index
is $2$. The formulas of
\citet[Proposition 6.1.2(i),(iii)]{A_Suzuki2024}
(Kato's symbol filtration) identify
\[
 \operatorname{gr}^0H^2(L,\F_2)=E^\times/E^{\times2},\qquad
 \operatorname{gr}^1H^2(L,\F_2)=\Omega^1_E,
\]
with $v\mapsto\{\widetilde v,2\}$ in degree zero and
$x\,d\log y\mapsto\{1+2\widetilde x,\widetilde y\}$ in degree one.
The restriction $H^2(R,\F_2)\hookrightarrow H^2(L,\F_2)$ is injective
and strict for this filtration; its degree-zero and degree-one images
are respectively
\[
 B^\times/B^{\times2}\subset E^\times/E^{\times2},\qquad
 \Omega^1_B\subset\Omega^1_E.
\]
These are exactly
\citet[Propositions 8.1.4, 8.2.1(i),(iii), and 8.2.3]{A_Suzuki2024}.
All their hypotheses hold: $A$ is regular, excellent and henselian,
$\zeta_2=-1\in A$, the unique prime dividing $2$ is $(2)$, and
$A/(2)=B$ is regular.

Suppose a class of the form \eqref{A:wild-image-form} belongs to
$U^1H^2(L,\F_2)$. We may write $u=2^e u_0$ with $u_0\in W^\times$.
Since $k$ is perfect, dividing $u_0$ by a square of a Teichm\"uller
unit puts it in $1+2W$ without changing its Kummer class.
The degree-zero image of $\{t,u\}$ is $t^e$ modulo squares.
The degree-zero image of $a$ is a unit of $B$ modulo squares.
Its $t$-adic valuation is even. Cancellation in degree zero therefore
forces $e$ to be even. Removing this square factor makes
$\{t,u\}\in U^1$, and strictness forces $a\in U^1$ as well.
Write $u=1+2b$ with $b\in W$. Its degree-one image is
$\bar b\,d\log t$ (the sign is immaterial in characteristic two).
Thus the degree-one image of every class in
\[
 \operatorname{im}\bigl(H^2(V,\F_2)\to H^2(L,\F_2)\bigr)\cap U^1
\]
belongs to
\begin{equation}\label{A:wild-form-obstruction}
 \Omega^1_B+k\,d\log t=B\,dt+k\,t^{-1}dt.
\end{equation}
On the other hand,
\[
 c=\{1+2/t,t\}\in U^1H^2(L,\F_2)
\]
has degree-one image $t^{-1}d\log t=t^{-2}dt$, which is not in
\eqref{A:wild-form-obstruction}. To see that no differential-form
qualification is hidden here, $\{t\}$ is a $2$-basis of both $B$ and
$E$, as $k$ is perfect. Hence their absolute differential modules are
$B\,dt$ and $E\,dt$. The coefficient of $t^{-2}$ is an explicit
linear functional separating the displayed form from the displayed
subspace. Thus $c$ is not in the image in \eqref{A:wild-cokernel}, and
its boundary is a nonzero class in $H^3(U,\F_2)$.
\end{proof}

\begin{remark}[Source and scope]
Gabber's 2015 IH\'ES lecture, at 33:17--33:24, pointed us toward a
possible extra degree for affine opens meeting the characteristic-$p$
locus. We used its automatic transcript only as a search lead.
The proof above uses the opened written symbol calculations and exact
localization sequences, not that transcript as a theorem.
No priority is claimed for this example or its derivation.
\end{remark}

\begin{theorem}[Proved under stated hypotheses; a uniform mixed-characteristic replacement]
\label{A:wild-affine-upper}
Let $A$ be an excellent henselian local domain of mixed characteristic
$(0,p)$, of dimension $d$, with residue field $k$. For every affine open
$U\subset\Spec A$,
\[
 \cd_p(U)\leq2d-1+d_p(k)
       \leq2d-1+\delta(k)+\cd_p(k).
\]
For $d=1$ this is the fraction-field bound. For $d=2$ the first bound
is attained by Theorem~\ref{A:wild-affine-counterexample}. If
$U\subset D(p)$ the sharper equality in
Proposition~\ref{A:wild-generic-open} still applies.
\end{theorem}
\begin{proof}
Write $s=d_p(k)$ and $r=\delta(k)$. Infinite $r$ makes the assertion
vacuous; assume $r<\infty$, so $r\leq s\leq r+1$.
The inequality $s\leq r+\cd_p(k)$ was proved in
Proposition~\ref{A:wild-generic-open}.
If $U$ contains the closed point it is $\Spec A$, whose cohomology is
that of $k$ by henselian invariance. We therefore assume that $U$
misses the closed point.

We induct on $d$, proving vanishing for all $p$-primary torsion sheaves.
For $d=1$, a nonempty such $U$ is the fraction field, and
Gabber--Orgogozo Th\'eor\`eme 1.2 gives $\cd_p(U)=1+s$.
For $d=2$, the closed subset $Z=U\cap V(p)$ is a finite set of closed
points. For $z\in Z$, its residue field is the fraction field of the
one-dimensional excellent henselian local domain $A/\mathfrak p_z$.
The equal-characteristic part of the same theorem gives
$d_p(\kappa(z))=1+s$.
Normalize the henselization of $\mathcal O_{U,z}$. It is a finite
product of henselian discrete valuation rings, with residue fields
finite over $\kappa(z)$. The generic fields therefore have
$p$-cohomological dimension $2+s$, by the discrete-valuation case and
finite-extension invariance of $d_p$.
The henselian local ring itself has $p$-cohomological dimension at most
one. Localization and henselian excision give
$H^q_z(U,F)=0$ for $q>3+s$. On $U[1/p]$ the bound is $2+s$ by
Gabber--Orgogozo Th\'eor\`eme 6.1. The localization sequence proves
$\cd_p(U)\leq3+s$.

Assume now $d\geq3$ and put $N=2d-1+s$.
We spell out the support argument, adapting the proof of
\citet[Expos\'e X, Th\'eor\`eme 4.1]{A_SGA4}.
By induction, an affine open in an integral closed subscheme
$\Spec A/\mathfrak q$ of characteristic zero and dimension $e<d$
has $p$-cohomological dimension at most $2e-1+s$.
The same estimate holds for a finite union of such closed subschemes
of dimensions at most $e$. Indeed, the exact sequence for a union of
two closed subsets reduces this assertion to their intersection,
whose characteristic-zero components have smaller dimension and
whose characteristic-$p$ part is affine with $p$-cohomological
dimension at most one. For $e=1$ those intersections lie at the closed
point of $\Spec A$, which $U$ misses; for $e\geq2$ the possible extra
degree from the intersection is at most $2\leq2e-1+s$.
Nilpotents do not change the \'etale topos.

It follows that any sheaf zero at the generic point of $U$ has
cohomological dimension at most $N-2$: approximate it by constructible
subsheaves and cover each support by finitely many integral closed
subschemes. Characteristic-zero components have dimension at most
$d-1$; the characteristic-$p$ part has bound one. The closed-union
exact sequence gives at most $\max(N-2,2)=N-2$ since $d\geq3$.

Let $i:\Spec K\to U$ be the generic inclusion and let $G$ be a finite
$p$-torsion Galois module over $K=\Frac A$.
At a characteristic-zero point $y\in U$ of height $h$, the stalks of
$R^b i_*G$ vanish for $b>h$, by
Theorem~\ref{A:strict-field} applied to the components of the strict
localization. At a characteristic-$p$ point $y$ of height $h$, we have
\[
 \delta(\kappa(y))=d-h+r.
\]
For this identity, apply
\citet[Lemmes 2.1.4--2.1.6 and the Cohen-structure calculation in \S4]{A_GO2008}
to the excellent henselian local domain $A/\mathfrak p_y$:
completion preserves the finite $p$-rank and the completed
fraction field has $p$-rank $\dim(A/\mathfrak p_y)+r$.
Excellence gives $\dim(A/\mathfrak p_y)=d-h$.
Strict henselization preserves $p$-rank. The fraction fields of its
components have cohomological dimension at most
$h+\delta(\kappa(y))=d+r$ by Gabber--Orgogozo.
Thus $R^b i_*G=0$ for $b>d+r$.

If $1\leq b\leq d-1$, every constructible subsheaf of $R^b i_*G$
is supported on the union of $U\cap V(p)$ and finitely many
characteristic-zero closed subschemes of dimensions at most $d-b$.
The closed-union sequence, the induction hypothesis and the affine
characteristic-$p$ bound give
\[
 \cd_p\bigl(U,R^b i_*G\bigr)
    \leq\max\{2(d-b)-1+s,\,2\}\leq N-b-1.
\]
Here the notation means the upper vanishing degree for that sheaf;
filtered-colimit compatibility removes constructibility.
If $d\leq b\leq d+r$, its support lies entirely in the closed affine
characteristic-$p$ subset, so its cohomology vanishes above degree one.
This is sufficient because
\[
 N-b-1\geq d-2+s-r\geq1.
\]

In the Leray spectral sequence
$H^a(U,R^b i_*G)\Rightarrow H^{a+b}(K,G)$, all possible incoming
differentials to $H^q(U,i_*G)$ for $q>N$ vanish by these inequalities.
The abutment vanishes too, since $\cd_p(K)=d+s\leq N$.
Hence $H^q(U,i_*G)=0$ for $q>N$.
Finally, for any constructible $p$-torsion sheaf $F$ on $U$, the kernel
and cokernel of $F\to i_*i^*F$ are zero at the generic point. Their
bound $N-2$ just proved and the two associated short exact sequences
give $H^q(U,F)=0$ for $q>N$. Coefficient d\'evissage and filtered
colimits finish the proof for all $p$-primary torsion sheaves.
\end{proof}
